\chapter{The Language}

This is the reference chapter. It describes wick v0.1 completely and
precisely: how source text is turned into tokens, what types exist, how
variables are scoped, how expressions are formed and in what order they bind,
what statements there are, how control flows, and how functions are declared
and called. Where the tutorial was content to gesture, this chapter states
rules. The grammar in Appendix~A is the formal companion to the prose here.

\section{Lexical structure}

A wick program is a sequence of Unicode bytes, read left to right and top to
bottom, that the lexer turns into a stream of tokens. Whitespace and newlines
separate tokens but are otherwise insignificant --- wick has no
newline-sensitive rules and no semicolons. A statement ends where the grammar
says it ends, not at a line break.

\paragraph{Comments.} A comment begins with `//` and runs to the end of the
line. There are no block comments.

\begin{lstlisting}[language=wick]
let x = 10   // this is a comment; it ends at the newline
\end{lstlisting}

\paragraph{Identifiers.} An identifier matches `[A-Za-z_][A-Za-z0-9_]*`: a
letter or underscore, then any number of letters, digits, or underscores.
Identifiers are case-sensitive. The following words are reserved and may not
be used as identifiers:

\begin{quote}
\ttfamily let\quad fn\quad if\quad elif\quad else\quad while\quad for\quad
in\quad break\quad continue\quad return\quad true\quad false\quad nil\quad
and\quad or\quad not
\end{quote}

\noindent
That is the entire keyword list --- seventeen words. Type names (`num`, `bool`,
`str`, `list`, `map`) are \emph{not} keywords; they are contextual, recognised
only where a type is expected, so nothing stops you from calling a variable
`str` if you are determined to confuse yourself.

\paragraph{Numbers.} There is one numeric type, `num`, an IEEE-754 double.
A numeric literal is a run of digits with an optional single decimal point:
`42`, `3.5`, `0.25`. There are no hexadecimal literals and no exponent
notation in v0.1. The lexer is careful about the range operator: `1..5` lexes
as the number `1`, the range token `..`, and the number `5` --- never as `1.`
followed by `.5`. If you want a literal with a fractional part next to a
range you must space it, but in practice ranges use integer bounds.

\paragraph{Strings.} A string literal is enclosed in double quotes and may
contain the escapes `\n` (newline), `\t` (tab), `\"` (a literal quote), and
`\\` (a literal backslash). Strings are immutable. There are no single-quoted
strings and no multi-line string literals.

\begin{lstlisting}[language=wick]
let a = "tenzin"
let b = "line one\nline two"
let c = "she said \"hi\""
\end{lstlisting}

\paragraph{Operators and punctuation.} The lexer recognises the following
multi-character tokens, greedily: `..` (range), `??` (nil-coalesce), `==`,
`!=`, `<=`, `>=`. The single-character tokens are the arithmetic and
comparison operators `+ - * / % < >`, the assignment `=`, the optional marker
`?`, the member dot `.`, and the grouping and separator characters
\texttt{( ) [ ] \{ \} , :}. Any other character is a lexical error,
reported as \texttt{unexpected character 'x'}.

\section{Types}

Every expression in wick has a static type, fixed at compile time. There are
six type constructors.

\begin{center}
\begin{tabular}{@{}llp{0.44\linewidth}@{}}
\toprule
\textbf{Type} & \textbf{Values} & \textbf{Notes} \\
\midrule
\texttt{num} & IEEE-754 double & the only number type \\
\texttt{bool} & \texttt{true}, \texttt{false} & the only type conditions accept \\
\texttt{str} & immutable text & \texttt{+} concatenates \texttt{str} with \texttt{str} only \\
\texttt{list<T>} & ordered, 0-based & \texttt{T} is \texttt{num}, \texttt{bool}, or \texttt{str} \\
\texttt{map<T>} & string-keyed & values are \texttt{num}, \texttt{bool}, or \texttt{str} \\
\texttt{T?} & \texttt{T} or \texttt{nil} & an \emph{optional} --- the headline feature \\
\bottomrule
\end{tabular}
\end{center}

The element type of a `list` or `map` must be one of the three scalar types
`num`, `bool`, or `str`. There are no nested containers in v0.1: `list<list<num>>`
and `map<list<str>>` do not exist, and attempting one produces
\texttt{container elements must be num, bool, or str}. Chapter~5 gives the
index-arithmetic workaround.

\paragraph{\texttt{nil} and optionals.} `nil` is not a value of general type.
It inhabits only optionals: an expression of type `T?` may be `nil` or may
hold a `T`, and there is no way to write a bare `nil` whose type is unknown.
`let x = nil` does not compile; the compiler cannot infer what the optional is
optional \emph{of}, and says so:
\texttt{can't infer a type from nil; annotate: let x: str? = nil}. Supply the
annotation and it is fine:

\begin{lstlisting}[language=wick]
let x: str? = nil     // x is str?, currently nil
\end{lstlisting}

\noindent
The whole of Chapter~4 is about optionals; here we note only that `T?` is a
first-class type that the compiler tracks, and that you cannot use a `T?`
where a `T` is required.

\paragraph{\texttt{void}.} A function declared without a return type has
return type `void`. `void` is not a value: a call to a void function is a
statement, not an expression, and cannot appear on the right of a `let` or
inside a larger expression. Attempting it yields
\texttt{can't assign a void expression}.

\section{Variables and scope}

Every variable is introduced by `let`. There is no other way to bring a name
into existence --- in particular, assigning to a name that was never declared
is not an implicit declaration but a compile error. This single rule removes
the entire class of typo-creates-a-global bug that haunts Lua.

\begin{lstlisting}[language=wick]
let x = 10          // num, inferred from 10
let y: num = 10     // the same, written out
let s = "hi"        // str
\end{lstlisting}

\noindent
The type may be inferred from the initialiser or written explicitly after a
colon. You must write it when there is nothing to infer from --- an empty
container literal or a bare `nil` --- and you may write it anywhere for clarity.
When both are present they must agree.

\paragraph{Global versus local.} A `let` at the top level of the file creates
a program \emph{global}: it is initialised once when the file loads and it
persists for the life of the program, across every frame. A `let` inside a
function creates a \emph{local}: it lives in that call's stack frame and is
gone when the call returns. This is why per-frame state --- the player
position, the score --- lives in top-level `let`s, while scratch values live in
locals.

\paragraph{Blocks and shadowing.} A variable is scoped to the block --- the
brace-delimited region --- in which it is declared, and to the blocks nested
inside it. An inner block may \emph{shadow} a name from an outer scope by
declaring a new `let` with the same name; the inner name hides the outer one
until the inner block ends. Redeclaring a name in the \emph{same} scope is an
error:

\begin{lstlisting}[language=wick]
let n = 1
if true {
  let n = 2      // ok: shadows the outer n inside this block
}
let n = 3        // ERROR: 'n' already declared in this scope
\end{lstlisting}

\noindent
The error for the last line is \texttt{'n' already declared}. To change an
existing variable you assign to it; you do not re-`let` it.

\paragraph{Assignment.} Assignment uses a single `=` and requires the name to
have been declared. Reading an undeclared name reports
\texttt{unknown variable 'x' (declare it with let)}; assigning to one reports
\texttt{unknown variable 'x' --- wick has no implicit globals; use let}. The
assigned value must have a type compatible with the variable (see the
compatibility rules below).

\section{Expressions and precedence}

Expressions combine literals, variables, calls, indexing, and operators.
Operators bind according to the following precedence table, from loosest
(evaluated last, binds least tightly) to tightest. All binary operators are
left-associative.

\begin{center}
\begin{tabular}{@{}cll@{}}
\toprule
\textbf{Level} & \textbf{Operators} & \textbf{Meaning} \\
\midrule
1 & \texttt{or} & logical or (short-circuit) \\
2 & \texttt{and} & logical and (short-circuit) \\
3 & \texttt{??} & nil-coalesce \\
4 & \texttt{==} \quad \texttt{!=} & equality, inequality \\
5 & \texttt{<} \enskip \texttt{<=} \enskip \texttt{>} \enskip \texttt{>=} & ordered comparison \\
6 & \texttt{+} \quad \texttt{-} & add / concatenate, subtract \\
7 & \texttt{*} \enskip \texttt{/} \enskip \texttt{\%} & multiply, divide, modulo \\
8 & unary \texttt{-}, \texttt{not} & negation, logical not \\
9 & \texttt{f(x)}, \texttt{a[i]} & call, index \\
\bottomrule
\end{tabular}
\end{center}

\noindent
So `a or b and c` parses as `a or (b and c)`, and `1 + 2 * 3` is `1 + (2*3)`,
as arithmetic convention demands. The one surprise for newcomers is that `??`
binds looser than comparison and arithmetic but tighter than `and`/`or`, which
is why `num(s) ?? 0` works without parentheses --- the coalesce applies to the
whole `num(s)`, and the result is a plain `num` ready for use.

\paragraph{Arithmetic.} `+ - * / %` operate on `num` and produce `num`. Both
operands must be non-optional `num`; using an optional is an error until you
unwrap it. Division by zero follows IEEE-754 (it produces infinity or NaN
rather than trapping), as in Lua.

\paragraph{The two faces of \texttt{+}.} `+` is numeric addition when both
operands are `num`, and string concatenation when both are `str`. It is never
mixed. There is no implicit conversion: `"score " + 5` does not compile,
because one side is `str` and the other `num`. The error is explicit about
the fix: \texttt{'+' needs num+num or str+str (use str(x))}. Write
`"score " + str(5)`. This is deliberate --- the silent stringification of
numbers is a rich source of bugs, and wick would rather you say what you mean.

\paragraph{Comparison and equality.} The ordered comparisons `< <= > >=`
require `num` operands and produce `bool`. Equality `==` and inequality `!=`
compare two values of the same type and produce `bool`; for `str`, equality
compares contents, not identity. Comparing values of different types is an
error (\texttt{'==' operands must have the same type}), with one blessed
exception: comparing an optional to `nil`, which is the narrowing idiom
covered in Chapter~4. Comparing a \emph{non}-optional to `nil` is itself an
error --- \texttt{comparing non-optional to nil} --- because a non-optional can
never be `nil`, so the check is dead code.

\paragraph{Logical operators.} `and` and `or` require `bool` operands, produce
`bool`, and short-circuit: `a and b` does not evaluate `b` when `a` is false,
and `a or b` does not evaluate `b` when `a` is true. `not` requires a `bool`
and negates it. Because there is no truthiness, you cannot write `a and b`
where `a` is a number; you write the comparison that yields the `bool` you
mean.

\paragraph{The nil-coalesce \texttt{??}.} `a ?? b` requires `a` to be an
optional `T?`; it evaluates to `a` unwrapped when `a` is non-`nil`, and to `b`
otherwise. The default `b` must have the base type `T`, and the whole
expression has type `T` --- non-optional. The right side is evaluated only when
the left is `nil`. Applying `??` to a non-optional left side is an error
(\texttt{'??' left side must be an optional}), because there is nothing to
coalesce.

\paragraph{Combined conditions and parentheses.} When a narrowing comparison
`x != nil` is combined with `and` or `or`, the compiler asks you to
parenthesise, so that the reader --- and the narrowing analysis --- can see the
structure unambiguously. The message is
\texttt{parenthesize this condition: (a != b) and ...}. Adding the parentheses
resolves it:

\begin{lstlisting}[language=wick]
if (s != nil) and ready { }    // parenthesised as the compiler asks
\end{lstlisting}

\section{Statements}

A wick program body is a sequence of statements. The statement forms are:

\begin{lstlisting}[language=wick]
let x = expr           // declaration (with optional : type)
x = expr               // assignment to a declared variable
xs[i] = expr           // element assignment (list index or map key)
f(a, b)                // call statement; a non-void result is discarded
if c { } elif c2 { } else { }
while c { }
for i in a..b { }      // i: num, from a to b-1
break                  // exit the innermost loop
continue               // next iteration of the innermost loop
return expr            // or bare `return` in a void function
\end{lstlisting}

\noindent
Blocks always use braces; there is no single-statement form without them. A
call used as a statement discards its result if it has one, which is how you
call a function for its effect. Element assignment works on a list index or a
map key; the target must already be a list or map (indexing anything else is
\texttt{only lists and maps can be indexed}).

\section{Control flow}

\paragraph{Conditionals.} `if`, optional `elif` chains, and an optional `else`
form the conditional. Every condition must be `bool`. There is no truthiness,
so `if 1 { }` and `if name { }` do not compile; the error names the rule
directly: \texttt{if condition must be bool (wick has no truthiness)}. Write
the comparison you mean --- `if x > 0`, `if s != ""`.

\begin{lstlisting}[language=wick]
if score > best {
  best = score
} elif score == best {
  // tie
} else {
  // nothing
}
\end{lstlisting}

\paragraph{\texttt{while}.} A `while` loop runs its body as long as the
`bool` condition holds. The condition is checked before each iteration.

\begin{lstlisting}[language=wick]
let i = 0
while i < 4 {
  i = i + 1
}
\end{lstlisting}

\paragraph{\texttt{for} ranges.} `for i in a..b` iterates the loop variable
`i` over the half-open range \texttt{[a, b)} --- from `a` up to but not
including `b`. Both bounds are `num` expressions, evaluated once. The loop
variable is a fresh local, scoped to the loop body, of type `num`. This is the
only iteration form in v0.1: there is no `for x in list`. To walk a list you
iterate its indices:

\begin{lstlisting}[language=wick]
for i in 0..len(xs) {
  lt.print(str(xs[i]), 4, 4 + i * 10, 1, 1, 1, 1)
}
\end{lstlisting}

\noindent
End-exclusivity is a deliberate match to 0-based indexing: `0..len(xs)` visits
exactly the valid indices `0` through `len(xs) - 1`, with no off-by-one to get
wrong.

\paragraph{\texttt{break} and \texttt{continue}.} `break` exits the innermost
enclosing loop; `continue` skips to that loop's next iteration. Both outside
any loop are errors: \texttt{break outside a loop},
\texttt{continue outside a loop}.

\section{Functions}

A function is declared with `fn`, at the top level only --- functions do not
nest, and there are no closures or first-class function values in v0.1.
Attempting a nested `fn` reports \texttt{fn declarations can't nest}.

\begin{lstlisting}[language=wick]
fn dist(ax: num, az: num, bx: num, bz: num): num {
  let dx = bx - ax
  let dz = bz - az
  return sqrt(dx * dx + dz * dz)
}
\end{lstlisting}

\paragraph{Parameters and return type.} Every parameter must be written with
its type; there is no inference for parameters. The return type, if any,
follows the parameter list after a colon. Omit it and the function is `void`.
A `void` function may use bare `return` to exit early but may not
`return expr` (\texttt{void function can't return a value}); a function with a
declared return type must `return expr` of that type on the paths that
produce a value, and may not `return` bare
(\texttt{this function must return T}).

\paragraph{Declare before use.} A function must be declared textually above
its first call. The compiler is single-pass and resolves each call as it
reaches it, so a call to a function defined later in the file reports
\texttt{unknown function 'f' (wick requires declare-before-use)}. The remedy
is to move the definition up. The two host-called functions, `update` and
`draw`, are exempt from ordering relative to each other, because the host --
not your code --- looks them up by name each frame.

\paragraph{Recursion.} A function may call itself, since its own name is in
scope within its body:

\begin{lstlisting}[language=wick]
fn fib(k: num): num {
  if k < 2 { return k }
  return fib(k - 1) + fib(k - 2)
}
\end{lstlisting}

\paragraph{Call checking.} Calls are fully type-checked. Passing the wrong
number of arguments reports \texttt{f() needs at least N arguments} or
\texttt{f() takes N arguments}; passing an argument of the wrong type reports,
for instance, \texttt{f() argument 3: expected num, got str}, naming the
function and the argument position. This applies equally to your own functions
and to engine calls, which are registered with typed signatures (Chapter~6).

\paragraph{Return-path coverage.} v0.1 does not yet verify that every path
through a non-void function returns a value. If control falls off the end of a
function that declares a return type, the call yields `nil` --- do not rely on
this. It is listed among the limits in Chapter~8, and the discipline is
simple: make the last statement of a value-returning function a `return`.

\section{Assignment compatibility}

The rules that govern when one type may be stored where another is expected
are short and worth stating together.

\begin{itemize}
\item A `T` may be stored where a `T?` is expected --- widening to an optional
is free and always safe.
\item A `T?` may \emph{not} be stored where a `T` is expected --- you must
unwrap it first, with `??` or by narrowing.
\item Containers are invariant: a `list<num>` accepts only `num` elements, and
a `list<num>` is not a `list<num?>` or anything else. The element type is part
of the type and must match exactly.
\item `void` cannot be stored at all, because it is not a value.
\end{itemize}

\noindent
These four rules, together with the optional discipline of the next chapter,
are the whole of wick's type checking. There is no subtyping beyond
`T`-into-`T?` widening, no coercion, and no generics beyond the built-in
`list` and `map`. The system is small on purpose --- small enough to hold in
your head, which is the point of the whole language.
